\documentclass[a4paper,12pt]{article}
\usepackage[utf8]{inputenc}
\usepackage[T1]{fontenc}
\usepackage[ngerman]{babel}
\usepackage{amsmath}
\usepackage{amssymb}
\usepackage{enumitem}
\usepackage{geometry}
\geometry{a4paper, margin=2.5cm}

\title{Einsendeaufgaben -- Aufgabe 3.2\\[0.5em]
\large 61211 Analysis}
\author{}
\date{}

\begin{document}
\maketitle

\section*{Aufgabe 1}

\begin{enumerate}[label=\alph*)]

    \item
    Für $(x,y) \neq (0,0)$ gilt $f(x,y) = \dfrac{x^2 y^2 (x-y)}{2x^2 + 2y^2}$.
    Folglich ist $f$ für $(x,y) \neq (0,0)$ wegen der Eigenschaft einer rationalen Funktion stetig.

    Noch $(x,y) = (0,0)$ zu untersuchen. Wir verwenden das $\varepsilon$-$\delta$-Kriterium.
    Sei $\varepsilon > 0$. Wir wählen $\delta := \min(1, \varepsilon)$.
    Sei $\binom{x}{y} \in \mathbb{R}^2 \setminus \{0\}$ mit $\left\|\binom{x}{y}\right\|_\infty < \delta$.

    Dann gilt:
    \begin{align*}
        |f(x,y)|
        &= \frac{x^2 y^2 \, |x-y|}{2(x^2+y^2)} \\
        &\leq \frac{x^2 (x^2+y^2) \, |x-y|}{2(x^2+y^2)} \\
        &= \frac{x^2 \, |x-y|}{2} \\
        &< \frac{|x-y|}{2} && \text{(wegen $|x| < \delta \leq 1$)} \\
        &\leq \frac{|x|+|y|}{2} \\
        &< \frac{2\varepsilon}{2} && \text{(wegen $\left\|\binom{x}{y}\right\|_\infty < \delta \leq \varepsilon$)} \\
        &= \varepsilon
    \end{align*}

    \item
    \[
        D_1 f(x,y) =
        \begin{cases}
            \dfrac{(3x^2 y^2 - 2xy^3)(2x^2+2y^2) + x^2 y^2 (x-y) \, 4x}{(2x^2+2y^2)^2}, & \text{falls } (x,y) \neq (0,0) \\[0.8em]
            0, & \text{falls } (x,y)=(0,0)
        \end{cases}
    \]

    \[
        D_2 f(x,y) =
        \begin{cases}
            \dfrac{(2yx^3 - 3y^2 x^2)(2x^2+2y^2) + x^2 y^2 (x-y) \, 4y}{(2x^2+2y^2)^2}, & \text{falls } (x,y) \neq (0,0) \\[0.8em]
            0, & \text{falls } (x,y)=(0,0)
        \end{cases}
    \]

    Da $D_1 f$ und $D_2 f$ für $\binom{x}{y} \neq \binom{0}{0}$ rationale Funktionen sind, sind sie stetig in $\binom{x}{y} \neq 0$.
    Die Stetigkeit von $D_1 f$ und $D_2 f$ in $\binom{x}{y} = 0$ ist noch zu untersuchen.

    Sei $\varepsilon > 0$ beliebig. Wir wählen $\delta := \min(1, \varepsilon)$.

    Dann gilt:
    \begin{align*}
        \frac{|3x^2 y^2 - 2xy^3|}{2x^2+2y^2}
        &\leq \frac{3x^2 y^2 + 2|xy^3|}{2x^2+2y^2} \\
        &\leq \frac{3x^2 + 2y^2}{2x^2+2y^2} && \text{(wegen $\left\|\binom{x}{y}\right\|_\infty < \delta \leq 1$)} \\
        &= x^2 \\
        &< \varepsilon && \text{(wegen $\left\|\binom{x}{y}\right\|_\infty < \delta \leq \varepsilon$ und $\left\|\binom{x}{y}\right\|^2 < \delta^2 \leq 1$)}
    \end{align*}

    Daraus folgt $\displaystyle\lim_{\binom{x}{y} \to 0} \frac{3x^2 y^2 - 2xy^3}{2x^2+2y^2} = 0$.

    Außerdem für $\delta := \min\!\left(1, \tfrac{\varepsilon}{4}\right)$ gilt
    \begin{align*}
        \frac{|x^2 y^2 (x-y) \, 4y|}{(2x^2+2y^2)^2}
        &\leq \frac{(2x^2+2y^2)^2 \, |x-y| \, 4|y|}{(2x^2+2y^2)^2} \\
        &= |x-y| \cdot 4|y| \\
        &\leq 4|x-y| && \text{(wegen $\left\|\binom{x}{y}\right\|_\infty < \delta \leq 1$)} \\
        &\leq 4|x| + 4|y| \\
        &< \varepsilon && \text{(wegen $\left\|\binom{x}{y}\right\|_\infty < \delta \leq \tfrac{\varepsilon}{4}$)}
    \end{align*}

    Daraus folgt
    \[
        \lim_{\binom{x}{y} \to 0} \frac{x^2 y^2 (x-y) \, 4y}{(2x^2+2y^2)^2} = 0.
    \]

    Damit können wir schließen:
    \begin{align*}
        \lim_{\binom{x}{y} \to 0} D_1 f(x,y)
        &= \lim_{\binom{x}{y} \to 0} \frac{2yx^3 - 3y^2 x^2}{2x^2+2y^2}
         + \lim_{\binom{x}{y} \to 0} \frac{x^2 y^2 (x-y) \, 4y}{(2x^2+2y^2)^2} \\
        &= 0 = D_1 f(0,0)
    \end{align*}

    $D_1 f$ ist also stetig in $0$.
    Mit der gleichen Argumentation ist $D_2 f$ stetig in $0$.
    Da also $D_1 f$ und $D_2 f$ stetig auf ganz $\mathbb{R}^2$ sind, ist $f$ nach Satz 3.6.4 total differenzierbar auf ganz $\mathbb{R}^2$.

    \item
    \begin{align*}
        D_v f(1,0)
        &= D_1 f(1,0) \cdot \frac{1}{\sqrt{2}} + D_2 f(1,0) \cdot \frac{1}{\sqrt{2}} \\
        &= \frac{1}{\sqrt{2}} (0 + 0) = 0
    \end{align*}
    \begin{align*}
        D_v f(0,0)
        &= D_1 f(0,0) \cdot \frac{1}{\sqrt{2}} + D_2 f(0,0) \cdot \frac{1}{\sqrt{2}} \\
        &= \frac{1}{\sqrt{2}} (0 + 0) = 0
    \end{align*}

\end{enumerate}

\end{document}
